\documentclass[11pt,a4paper]{article}
\usepackage[margin=23mm]{geometry}
\usepackage{fontspec,xeCJK,amsmath,booktabs,tabularx,xcolor,fancyhdr,hyperref}
\setmainfont{Times New Roman}
\setCJKmainfont{Nanum Myeongjo}
\setCJKsansfont{Nanum Gothic}
\xeCJKsetup{CJKspace=true}
\definecolor{accent}{HTML}{155E63}
\hypersetup{colorlinks=true,urlcolor=accent,pdftitle={연속 판정 문턱에서의 부분 관측 인증},pdfauthor={김태건, 정지윤, 정지훈}}
\pagestyle{fancy}
\fancyhf{}
\lhead{\small\sffamily 강원과학고 R\&E}
\rhead{\small\sffamily AI·정보과학}
\rfoot{\small\thepage}
\setlength{\headheight}{15pt}
\setlength{\parindent}{1em}
\setlength{\parskip}{.45em}
\linespread{1.23}
\renewcommand{\arraystretch}{1.26}
\newcommand{\lead}[1]{\textcolor{accent}{\sffamily\bfseries #1}}
\begin{document}
\begin{center}
{\sffamily\small 연구보고서 \quad 2026년 11월 3일}\par\vspace{1.1em}
{\sffamily\bfseries\LARGE 연속 판정 문턱에서의 부분 관측 인증}\par\vspace{.55em}
{\sffamily\large 활–현 경계의 법칙 판정과 적응형 파형 조회}\par\vspace{1em}
김태건 · 정지윤 · 정지훈\par
강원과학고등학교
\end{center}
\section*{초록}
자동 판정 규칙이 달라지면 신호에서 추정한 경계와 그 경계로 내린 결론도 달라질 수 있다. 본 연구는 공개 첼로 로봇 활 자료에서 두 공통 문턱의 연속 영역을 대상으로, 일부 파형만 확인한 상태에서도 지정한 활 힘 법칙과의 양립 여부를 확정하는 알고리즘을 개발하였다. 조회한 특징으로 문턱 영역을 분할하고, 미조회 파형이 만들 수 있는 첫 양성 연속쌍을 절편 구간에 전파한다. 이 구간을 이용해 한 영역을 배제하거나 양립하는 규칙의 존재를 인증한다. 이미 분석에 사용한 12개 속도 폴더에서 두 허용 거리를 적용하여 양립 2건과 불양립 22건을 얻었다. 세 조회 정책의 판정 72건은 전체 파형의 특징을 이용한 전수 계산과 일치하였다. 고정 9규칙 실험에서 공동 응답이 절편 교집합에 미치는 변화를 평가하는 정책은 허용 거리 0.2의 조회를 1,616회에서 910회로 줄였다. 연속 인증에서 학습 확률의 추가 절감은 허용 거리 0.2의 한 폴더에서 관찰되었다. 주된 기여는 연속 문턱, 첫 경계, 법칙 판정을 부분 관측의 종료 조건으로 연결한 계산 방법이다.

\noindent\lead{주제어}\quad 능동 조회, 제약 추론, 부분 관측, 연속 문턱, 활–현 경계

\section{연구 문제와 자료}
같은 파형에 다른 문턱을 적용하면 헬름홀츠 운동을 가리키는 자동 라벨이 바뀐다. 이때 힘 순서에서 처음 나타나는 양성 쌍이 이동하고, 여러 활 위치에서 관측한 경계가 같은 식과 양립하는지도 달라진다. 본 연구는 지정한 영역에 양립하는 공통 문턱이 존재하는지를 판단하고, 판정에 충분한 증거를 얻으면 조회를 종료한다.

빈 음대 IWK의 공개 로봇 첼로 G현 자료를 사용하였다 [1]. 한 측정 열은 같은 속도 폴더 안의 한 활 위치를 뜻한다. 각 열에는 힘으로 정렬한 서로 다른 긋기 50개가 있다. 한 번 조회하면 해당 파형의 특징과 판정 조건을 공개한다. 힘과 활 위치는 미리 주어진 좌표로 취급하였다. 대상 열 목록은 기본 문턱에서 유한한 양의 경계를 가진 열로 미리 정했다. 조회 비용은 이 목록이 주어진 상황에 조건부이다.

\clearpage
\section{연속 문턱 인증}
주기성 점수를 \(P\), 둘째 되돌이 비율을 \(R\), 물리적 통과 조건을 \(g\), 문턱과 무관하게 양성인 조건을 \(f\)라고 하자. 파형 라벨을 다음과 같이 정의한다.
\begin{equation}
H_{u,v}=f\lor(g\land[P\ge u]\land[R<v]),
\qquad (u,v)\in[0.7,0.9]\times[0.3,0.7].
\end{equation}
두 문턱은 모든 대상 열에 공통이다. 각 열에서 힘순 첫 양성 연속쌍의 시작 위치를 찾는다. 시작 위치 \(i\)의 경계는 \([F_{i-1},F_i]\)이다. 연속쌍이 없거나 시작이 첫 관측이거나 양의 폭을 갖는 양의 힘 구간이 아니면 해당 후보를 무효로 처리한다.

활 위치 \(\beta_c\)의 경계가 \(\log F=a+p\log\beta_c\)와 양립하려면 같은 절편 \(a\)가 존재해야 한다. \(p=-2\)로 고정하고, 허용 로그 힘 거리 \(\varepsilon\)를 양 끝에 적용하였다. 열 \(c\)의 유효 후보 \(i\)가 허용하는 절편 구간은
\begin{equation}
I_{ci}=[\log F_{i-1}-p\log\beta_c-\varepsilon,\
          \log F_i-p\log\beta_c+\varepsilon].
\end{equation}
완전 관측에서 모든 대상 열의 경계가 유효하며 절편 구간에 공통점이 있으면 그 규칙이 양립한다.

\subsection*{관측으로부터 문턱 영역 구성}
조회한 파형의 \(P,R\)에서만 절단점을 만든다. 사각형에서 알려진 라벨의 최솟값은 \(u\)가 높고 \(v\)가 낮은 모서리에서, 최댓값은 반대 모서리에서 얻는다. 두 값이 같으면 사각형 전체에서 라벨을 고정한다. 값이 다른 파형과 미조회 파형은 미정으로 둔다. 각 축의 셀은 오른쪽 끝점을 포함하며, 맨 왼쪽 끝점도 포함한다. 이는 \(P\ge u\)와 \(R<v\)의 등호 처리를 보존한다.

특정 위치가 첫 쌍이 되려면 그 위치와 다음 위치를 모두 1로 둘 수 있어야 한다. 앞에서 이미 확정된 쌍이 없어야 하고, 직전 위치가 있다면 0으로 둘 수 있어야 한다. 나머지 미정 비트를 0으로 두면 이 조건을 만족하는 첫 쌍을 구성할 수 있다. 확정된 쌍이 없으면 쌍이 없는 완성도 후보에 포함한다.

\subsection*{종료 조건}
각 열에서 모든 유효 후보의 절편 구간을 감싸는 외부 구간을 구한다. 유효 후보가 없는 열이 있거나 외부 구간끼리 공통점이 없으면 사각형 전체를 배제한다. 모든 후보가 유효하고 모든 열의 모든 후보 구간이 공통 절편을 가지면 양립을 확정한다. 후보를 계산할 때 미정 비트를 자유롭게 둘 수 있도록 완화하므로 실제 가능한 완성은 모두 포함된다. 이 완화는 조기 판정을 늦출 수 있다.

판정이 남으면 알려진 절단점에서 분할한다. 양립을 보장하는 영역 하나로 존재 판정을 끝내며, 불양립은 전 영역을 배제한 뒤 확정한다. 모든 파형을 조회하면 가장 작은 셀에서 각 라벨이 고정돼 전수 계산과 일치한다.

\clearpage
\section{AI 조회 정책과 실험 설계}
\noindent\fbox{\parbox{.95\linewidth}{\centering\sffamily
조회된 특징 \(\longrightarrow\) 가능한 경계와 절편 \(\longrightarrow\) 종료 인증\\[4pt]
미결이면 다음 파형 선택 \(\longrightarrow\) 특징 공개 \(\longrightarrow\) 인증 갱신}}
\par\medskip
인증기는 지금까지 관측한 정보로 양립 여부를 확정할 수 있는지 판단한다. 미결이면 조회 정책은 각 파형의 가능한 응답에 따라 다음 관측의 가치를 계산한다. 조회 정책을 교체할 때도 같은 종료 기준으로 비교할 수 있다.

\subsection*{공동 응답과 절편 변화에 기반한 선택}
조회 정책에서는 계산량을 제한하기 위해 고정된 문턱 쌍 9개를 앵커로 사용하였다. 파형 하나의 주기성과 비율이 9개 라벨을 함께 결정한다. 가능한 공동 응답은 양성 영역을 갖는 직사각형 패턴 9개와 전부 음성인 패턴을 합쳐 10개다.

후보 파형마다 가능한 공동 응답을 가정한다. 바뀌는 열 하나의 첫 경계와 절편 구간을 다시 계산하고, 다른 열의 저장된 절편 극값과 합친다. 예상 생존 규칙 수가 가장 적은 파형을 고르고, 동률이면 남은 교집합 폭의 합이 작은 파형을 고른다. 앵커의 판정이 정리돼도 연속 영역이 미결이면 조회를 계속한다.

\subsection*{학습 확률의 추가 효과를 분리하는 대조}
공동 응답에 균등 확률을 주는 정책과 공개 힘·위치·속도로 공동 응답의 확률을 추정하는 HGB 정책을 비교하였다. HGB는 A/B 측정 계열 1,080개 열에서 한 번 적합하였다. C/D 평가 파형의 특징은 해당 파형을 조회할 때 선택 상태에 공개한다. 힘 순위 우선도 함께 비교하였다. 세 정책은 동일한 연속 인증기를 사용하였다.

\subsection*{평가 규약과 구현 확인}
연속 실험은 이미 사용한 C/D 계열 12개 속도 폴더에 허용 거리 0.2와 0.4를 적용하였다. 같은 폴더에서 나온 두 결과는 서로 의존한다. 각 정책은 시작 시점, 이후 64회마다, 모든 파형을 읽은 시점에 인증을 검사했다. 부분 검사에서 사각형 192개 한도를 만나면 미결 상태로 조회를 이어 갔다. 비용은 저장 특징을 이용한 모의 조회 횟수로 집계하였다.

완전 관측 결과는 기존 연속 영역 전수 계산과 대조하였다. 별도로 길이 6의 합성 두 열 100개 사례에서 셀 직접 열거와 비교했다. 합성 양립과 불양립에서 모두 완전 관측 결과가 일치했으며, 중간에 확정한 판정도 완전 관측 판정과 일치했다.

\clearpage
\section{결과}
\subsection*{연속 문턱의 조기 인증}
허용 거리 0.2에서 실제 12개 폴더 모두 불양립이었다. 0.4에서는 양립 2개와 불양립 10개였다. 이 24개 폴더·허용 거리 조합에 세 정책을 적용한 판정 72건은 모두 전수 판정과 일치했다.

\begin{center}
\small
\begin{tabularx}{\linewidth}{Xrrr}
\toprule
평가 범위 & 힘 순위 우선 & 균등 응답 & 학습 응답 \\
\midrule
0.2 불양립 12건 & 1,856 & 1,280 & 1,216 \\
0.4 양립 2건 & 2,112 & 1,088 & 1,088 \\
0.4 불양립 10건 & 2,880 & 2,304 & 2,304 \\
\bottomrule
\end{tabularx}
\end{center}
\noindent 표 1. 연속 인증이 끝날 때까지의 총 모의 조회 수.

0.2에서 학습 응답은 균등 응답보다 64회 적었다. D2 세 번째 속도 폴더 한 곳에서 인증 검사가 한 구간 앞당겨진 결과다. 0.4에서는 두 응답 정책의 조회 수가 모든 폴더에서 같았다. 양립 두 폴더는 같은 C\_T2 세션에 속한다. 0.4에서 선택과 인증을 합친 실행 시간은 힘 순위 우선 0.87초, 균등 응답 4.82초, 학습 응답 4.34초였다. 실제 파형 처리 비용은 별도로 측정하지 않았다.

\subsection*{고정 9규칙에서 분리한 조회 정책 효과}
별도 실험은 고정 9개 규칙과 매 조회 검사 조건을 사용하였다. 같은 인증기를 적용했을 때 기존 법칙 판정 우선 정책의 총 1,616회가 공동 응답을 직접 평가하는 정책에서 910회로 줄었다. 허용 거리 0.2에서 \(43.69\%\)의 절감이다. 조회가 필요했던 네 폴더 중 세 곳에서 절감했고 D2에서는 36회 늘었다. 0.4의 양립 두 폴더에서는 1,070회에서 1,033회로 줄었다. 응답 가중치를 훈련 자료의 발생 빈도에서 균등 확률로 바꾸어도 조회 수가 같았다.

앞선 HGB 정책은 국소 경계 후보의 예상 제거 수를 점수로 썼으며, 0.2에서 1,725회 조회했다. D1과 D2에서 규칙 배제가 늦어져 다른 두 폴더의 절감을 넘는 추가 조회가 발생했다. 국소 후보를 없애도 절편 극값이 유지되면 양립 여부가 바뀌지 않는다. 공동 응답을 직접 평가하는 점수는 절편 교집합의 변화에 반응한다.

\subsection*{D2의 종료 병목}
고정 9규칙의 D2에서 새 정책은 25회와 202회에 각각 규칙 세 개를 배제하였다. 비율 문턱 0.5의 마지막 세 규칙은 339회까지 남았다. 기존 정책은 291회에 처음 배제했지만 303회에 모든 규칙을 정리하였다. 새 정책의 외부 교집합 폭 합은 256회와 303회에서 모두 약 1.472였다. 이 사례에서는 마지막 규칙을 판별할 증거가 조회 비용을 좌우했다.

\clearpage
\section{기여와 결론}
\lead{주 기여}\quad
관측으로 드러난 특징값으로 연속 문턱 영역을 구성하고, 미조회 파형이 만들 수 있는 첫 경계를 절편 구간으로 전파하였다. 이 연결을 통해 부분 관측에서도 영역 전체를 배제하거나 양립하는 규칙의 존재를 인증할 수 있다. 모든 파형을 확인한 상태의 공통 규칙 검증을 일부 파형만 확인한 상태의 종료 판정으로 확장하였다.

\lead{AI 기여}\quad
파형 하나가 줄 수 있는 공동 응답을 가정하고, 각 응답이 전체 법칙의 절편 교집합에 미치는 변화로 다음 파형의 우선순위를 정한다. 한 파형을 가정할 때 바뀌는 열만 다시 계산한다. 능동 의사결정에서 관측의 가치를 평가하는 접근 [2,3]과 최소 불만족 부분집합을 찾는 제약 설명 [4]에는 선행연구가 있다. 본 연구가 검증한 차별점은 공통 문턱, 힘순 첫 연속쌍, 법칙 절편의 결합을 부분 조회의 인증과 우선순위에 사용한 구체적인 구성이다.

\lead{결과의 범위}\quad
가장 직접적인 근거는 연속 영역에서의 판정 일치와 고정 9규칙에서 동일 인증기로 측정한 조회 절감이다. 학습의 추가 절감은 한 검사 구간에서 관찰되었다. 평가 자료는 이미 개발에 사용됐고 실제 양립 폴더는 한 측정 세션에 속한다. 독립 세션 검증과 대상 열 목록을 만드는 비용의 측정이 남아 있다. 판정의 대상은 자동 라벨에서 얻은 경계와 지정한 식의 양립성이다. 물리 법칙의 통계 검정이나 사람의 청취 판단과는 해석 범위가 다르다.

\section*{참고 자료}
\begin{enumerate}
\item A. Lampis. \textit{A collection of bowed responses of cello strings}. IWK, University of Music and Performing Arts Vienna. \href{https://zenodo.org/records/17749111}{Zenodo 컬렉션}.
\item S. Javdani et al. \textit{Near Optimal Bayesian Active Learning for Decision Making}. AISTATS, 2014. \href{https://proceedings.mlr.press/v33/javdani14.html}{PMLR}.
\item Y. Chen et al. \textit{Submodular Surrogates for Value of Information}. AAAI, 2015. \href{https://ojs.aaai.org/index.php/AAAI/article/view/9694}{논문}.
\item E. Gamba, B. Bogaerts and T. Guns. \textit{Efficiently Explaining CSPs with Unsatisfiable Subset Optimization}. IJCAI, 2021. \href{https://ijcai.org/proceedings/2021/191}{논문}.
\end{enumerate}

\noindent\lead{재현 근거}\quad
저장소 \texttt{law\_continuous\_v1}의 계획, 인증 코드, 기능시험, 폴더별 결과와 D2 진단을 따른다. \texttt{law\_counterfactual\_selection\_v1}에 고정 9규칙 비교를 보존하였다. 두 실험의 조회 수는 규칙 영역과 검사 간격이 다르므로 각 실험 안에서 비교한다.
\end{document}
